\honest{The ``Intuitions'' Behind ``Utilitarianism''}{Eliezer Yudkowsky}{2008}

I keep moral discussion at the level of cases, because metaethics depends on the Mind Projection Fallacy, which I have not yet discussed. \nb{In the book, both that post and the metaethics posts come before this one.}

Paul Gowder says that my choice of one person's torture over a googolplex dust specks, and my utilitarianism, both rest on intuition, and that I never argued for the utilitarian intuitions. \nb{Gowder's point was that an argument showing two beliefs inconsistent does not say which to give up, since both rest on intuitions ``on equal footing''.} I agree that intuitions are all there is: delete them and you are left with a rock. The work is ``renormalizing intuition'', keeping the ones that survive reflection. \nb{Moral philosophers call a method of this kind reflective equilibrium.} Refuse to build on the reflective ones and you are left running in circles, with circular preferences. \nb{The argument about circular preferences was in ``Circular Altruism''; the book's version of that post, ``Feeling Moral'', leaves it out.}

Gowder does not say what he means by ``utilitarianism'', so I list seven claims and say which I accept. I accept that two separate occurrences of a harm are exactly twice as bad as one.

Our intuitions about how to reach our goals are ``frankly messed up''. People give more to save one sick child than to save a group of eight sick children. The reason is obvious, I say: one child arouses more emotion than eight. \nb{The researchers, Kogut and Ritov, reached the same explanation. In a second study, when people saw the one child and the group side by side, most chose the group.} The intuition is wrong: ``the brain doesn't goddamn multiply.'' Five efforts that each save 5,000 lives do worse than one that saves 50,000. ``Three lives are one and one and one,'' however many others exist. That is aggregation. \nb{Many who refuse to prefer the torture accept this. Many contractualists, who follow Scanlon, hold that a rescuer should save five rather than one. What they deny is that small harms to many can outweigh a grave harm to one. Parfit wrote that we might spare one person a hundred days of pain rather than spare ``any number'' of others ten days. The argument from counting lives does not reach that case.}

People say no money is worth a life and then drive an extra mile to save \$10. Their absolute rules are social signals: a stated trade-off ``would sound like you were plotting to defect.'' They do not show two tiers of value. A lower tier would never decide anything, since some upper-tier consideration is always present. \nb{This objection is a strong one, and Michael Huemer later developed it against lexical priority in general. A version of it is pressed against contractualism, whose defenders have replied.}

The value of one event, such as music, may be complex. But ``When lives are at stake, I shut up and multiply.'' That is why I am a utilitarian, at least ``most of the time''.
